\documentclass[11pt]{amsart}
\usepackage{amsmath,amssymb,amsthm}
\usepackage[margin=1.2in]{geometry}
\usepackage{hyperref}

\newtheorem{theorem}{Theorem}[section]
\newtheorem{lemma}[theorem]{Lemma}
\newtheorem{corollary}[theorem]{Corollary}
\newtheorem{proposition}[theorem]{Proposition}
\theoremstyle{remark}
\newtheorem{remark}[theorem]{Remark}

\newcommand{\Cp}{\mathbb{C}_{+}}
\newcommand{\R}{\mathbb{R}}
\newcommand{\N}{\mathbb{N}}
\newcommand{\HH}{H^{2}(\Cp)}
\newcommand{\Hinf}{H^{\infty}(\Cp)}
\newcommand{\KT}{K_{\Theta}}
\newcommand{\dist}{\operatorname{dist}}
\newcommand{\re}{\operatorname{Re}}
\newcommand{\im}{\operatorname{Im}}

\title{The scalar zeta scattering semigroup is not uniformly\\
exponentially stable, unconditionally}
\author{OpenAI Codex}
\date{19 July 2026}

\begin{document}
\maketitle

\begin{abstract}
Let $\xi$ denote the completed Riemann zeta function, let
$\Theta(s)=\xi(2s)/\xi(1+2s)$, which is inner on the right half-plane, and
let $Z(t)$ be the truncated translation semigroup on the model space
$\KT = H^{2}\ominus\Theta H^{2}$. We prove unconditionally that
$\|Z(t)\|=1$ for every $t\ge 0$, so the growth bound of $Z$ is $0$;
equivalently, $\dist_{H^{\infty}}(e^{-ts},\Theta H^{\infty})=1$ for every
$t>0$. Under the Riemann hypothesis the point spectrum of the generator
lies on the line $\re\lambda=-\tfrac14$, so the spectral bound--growth
bound gap of this semigroup is maximal \emph{even assuming RH}. The
mechanism is arithmetic: by Selberg's theorem a positive proportion of
zeta zeros lie on the critical line, and these zeros cluster with
unbounded multiplicity in shrinking windows, producing explicit
normalized reproducing kernels on which the semigroup is asymptotically
norm-preserving. Consequently, no operator-norm form of the
Lax--Phillips energy-decay criterion for RH can hold for any realization
of the scattering semigroup that is norm-equivalent to this scalar
model; the per-vector (local energy decay) form of the criterion is
untouched.
\end{abstract}

\section{Introduction}

Lax and Phillips \cite{LaxPhillips} and Faddeev--Pavlov \cite{FaddeevPavlov}
connected the Riemann hypothesis to the decay of the energy semigroup
$Z(t)$ of the automorphic wave equation on the modular surface: RH is
equivalent to exponential decay of $Z(t)f$ for initial data $f$ of
compact support, at rate governed by the real parts of the resonances,
which sit at distance $\tfrac14$ from the imaginary axis if and only if
RH holds. It is tempting to formulate the criterion in operator norm,
\begin{equation}\label{eq:naive}
  \lim_{t\to\infty}\frac1t\log\|Z(t)\| \le -\tfrac14
  \qquad\text{(purportedly equivalent to RH),}
\end{equation}
and statements of this shape appear in the literature (e.g.\ Lax's
equation~(54) in \cite{LaxTodo}).

The purpose of this note is to show that \eqref{eq:naive} fails
\emph{unconditionally} for the canonical scalar realization of the
scattering semigroup: the truncated translation semigroup on the model
space of the inner function $\Theta(s)=\xi(2s)/\xi(1+2s)$, whose zeros
are the points $\rho/2$ over the nontrivial zeros $\rho$ of $\zeta$.
Precisely, $\|Z(t)\|=1$ for all $t$, so the uniform growth bound is $0$
regardless of the truth of RH, while the spectral bound is $-\tfrac14$
under RH. The proof is elementary once one observes that the on-line
zeros guaranteed by Selberg's theorem \cite{Selberg} cluster with
unbounded multiplicity: clusters of $m$ zeros make $|\Theta|$ decay like
$q^{m}$ at nearby points of small real part, and normalized reproducing
kernels at such points are asymptotic eigenvectors of $Z(t)^{*}$ with
eigenvalue modulus arbitrarily close to $1$.

Since the argument uses only zeros \emph{on} the critical line---which
exist in both the RH-true and RH-false worlds---it cannot bear on RH
itself. Its content is a correction of formulation: the RH-equivalent
statement in the Lax--Phillips criterion is the per-vector decay of
$Z(t)f$ for fixed admissible data, and any operator-norm strengthening
of it is false. In semigroup-theoretic language, the scalar zeta
scattering semigroup exhibits the maximal possible gap between spectral
bound and growth bound (cf.\ \cite{EngelNagel, GearhartPruess}), with
the gap forced by the arithmetic of the zero distribution. The
phenomenon fits the general framework of spectral singularities of
contractive semigroups of Naboko--Romanov \cite{NabokoRomanov}, here
with the singularity located at infinity.

\subsection*{Relation to the Lax--Phillips semigroup}
Our theorem concerns the scalar model semigroup on $\KT$ and is
self-contained. Transferring the conclusion to the energy-space
semigroup of \cite{LaxPhillips} requires a unitary (or norm-)
equivalence between that semigroup and the scalar compression; such an
identification is developed by Uetake \cite{Uetake}. We do not re-derive
it here, and we state the consequence for \eqref{eq:naive}
conditionally on it (Corollary~\ref{cor:LP}). Note that the full
scattering matrix $S_{c}(s)=-\frac{s-1/2}{s+1/2}\,\Theta(s)$ differs
from $\Theta$ by a single Blaschke factor, which changes neither the
cluster structure of the zero set nor, therefore, the conclusion of the
main theorem (Remark~\ref{rem:extra-factor}).

\section{Preliminaries}

Let $\Cp=\{s\in\mathbb C:\re s>0\}$ and let $\HH$, $\Hinf$ be the Hardy
spaces of the half-plane. For $\mu\in\Cp$ the reproducing kernel of
$\HH$ is
\[
  k_{\mu}(s)=\frac{1}{2\pi}\,\frac{1}{s+\bar\mu},
  \qquad
  \langle f,k_{\mu}\rangle=f(\mu),
  \qquad
  \|k_{\mu}\|^{2}=k_{\mu}(\mu)=\frac{1}{4\pi\,\re\mu}.
\]
For an inner function $\Theta$ on $\Cp$, the model space is
$\KT=\HH\ominus\Theta\HH$, with reproducing kernel
\[
  k^{\Theta}_{\mu}
  = P_{\KT}k_{\mu}
  = k_{\mu}-\overline{\Theta(\mu)}\,\Theta k_{\mu},
  \qquad
  \|k^{\Theta}_{\mu}\|^{2}
  = \bigl(1-|\Theta(\mu)|^{2}\bigr)\|k_{\mu}\|^{2}.
\]
For $\varphi\in\Hinf$ the truncated Toeplitz operator is
$A_{\varphi}=P_{\KT}M_{\varphi}|_{\KT}$. We will use two standard facts
(see \cite{GMR, Sarason}):
\begin{enumerate}
\item[(i)] $A_{\varphi}^{*}g = P_{H^{2}}(\bar\varphi g)$ for
  $g\in\KT$, and $K_{\Theta}$ is invariant under the co-analytic
  Toeplitz operator $T_{\bar\varphi}$; moreover
  $T_{\bar\varphi}k_{\mu}=\overline{\varphi(\mu)}\,k_{\mu}$ for every
  $\mu\in\Cp$.
\item[(ii)] (Sarason \cite{Sarason})
  $\|A_{\varphi}\|=\dist_{H^{\infty}}(\varphi,\Theta\Hinf)$.
\end{enumerate}
The \emph{truncated translation semigroup} on $\KT$ is
\[
  Z(t) := A_{\varphi_{t}},
  \qquad
  \varphi_{t}(s)=e^{-ts}
  \quad (t\ge 0),
\]
a strongly continuous contraction semigroup whose generator we denote
$A$; its point spectrum is $\{-\bar a:\Theta(a)=0\}$, with eigenvectors
of $Z(t)^{*}$ the kernels $k_{a}=k^{\Theta}_{a}$ at zeros $a$ of
$\Theta$. (Adjoint and sign conventions differ across the literature;
none affects operator norms, which is all we use.)

\section{The inner function}

\begin{lemma}\label{lem:inner}
Let $\xi(w)=\tfrac12 w(w-1)\pi^{-w/2}\Gamma(w/2)\zeta(w)$ and
\[
  \Theta(s)=\frac{\xi(2s)}{\xi(1+2s)}
  =\sqrt{\pi}\;\frac{2s-1}{2s+1}\,
   \frac{\Gamma(s)}{\Gamma(s+\tfrac12)}\,
   \frac{\zeta(2s)}{\zeta(2s+1)}.
\]
Then $\Theta$ is inner on $\Cp$. Its zero set in $\Cp$ is exactly
$\{\rho/2:\zeta(\rho)=0,\ 0<\re\rho<1\}$, with multiplicities, and there
is no cancellation: $\xi(1+\rho)\neq0$ for every nontrivial zero
$\rho$.
\end{lemma}

\begin{proof}
By the functional equation $\xi(w)=\xi(1-w)$ we have
$\xi(1+2s)=\xi(-2s)$, so the zeros of the denominator in $\Cp$ would
require $1+2s=\rho$, i.e.\ $\re s=(\re\rho-1)/2\in(-\tfrac12,0)$;
hence $\Theta$ is analytic on the closed right half-plane (the trivial
zeros of $\zeta(2s)$ lie at $\re s<0$ and are absorbed by $\Gamma$
poles in $\xi$; the apparent zero of $2s-1$ at $s=\tfrac12$ cancels the
pole of $\zeta(2s)$). Zeros in $\Cp$: $\xi(2s)=0$ iff $2s=\rho$, giving
$s=\rho/2$ with $\re(\rho/2)\in(0,\tfrac12)$; and
$\xi(1+\rho)\neq0$ since $\re(1+\rho)>1$.

Boundary modulus: since $\xi(\bar w)=\overline{\xi(w)}$ and
$\xi(1+2i\tau)=\xi(-2i\tau)=\overline{\xi(2i\tau)}$, we get
$|\Theta(i\tau)|=1$ for all $\tau\in\R$ ($\xi$ does not vanish on
$\re w=0$).

Growth: for $\re s\ge 2$ the factor $\zeta(2s)/\zeta(2s+1)$ is bounded
(absolute convergence) and $\Gamma(s)/\Gamma(s+\tfrac12)=O(|s|^{-1/2})$
by Stirling, uniformly in $|\arg s|\le\pi/2$; for $0<\re s<2$ one has
the finite-order bound $\zeta(2s)\ll(1+|s|)^{A}$ and the classical
lower bound $1/\zeta(\sigma+it)\ll(\log(|t|+2))^{7}$ for $\sigma\ge1$
\cite[Ch.~3]{Titchmarsh}, so $|\Theta(s)|\ll(1+|s|)^{B}$ on
$\overline{\Cp}$. A function analytic on $\Cp$, continuous on the
closure, of polynomial growth, with modulus $\le1$ on the boundary is
bounded by $1$ on $\Cp$ by the Phragm\'en--Lindel\"of principle for a
half-plane. Hence $\Theta\in\Hinf$ with unimodular boundary values:
$\Theta$ is inner.
\end{proof}

\section{Clusters of critical zeros}

\begin{lemma}\label{lem:clusters}
There exist $y_{n}\to\infty$, integers $m_{n}\to\infty$, and real
numbers $\varepsilon_{n,1},\dots,\varepsilon_{n,m_{n}}$ with
$\max_{j}|\varepsilon_{n,j}|\to0$, such that the points
\[
  a_{n,j}=\tfrac14+i\,(y_{n}+\varepsilon_{n,j}),
  \qquad j=1,\dots,m_{n},
\]
are distinct zeros of $\Theta$.
\end{lemma}

\begin{proof}
By Selberg's theorem \cite{Selberg} the number $N_{0}(T)$ of zeros of
$\zeta$ on the critical line with ordinate in $(0,T]$ satisfies
$N_{0}(T)\ge cT\log T$ for some $c>0$ and all large $T$. Cover $(0,T]$
by $\lceil T(\log T)^{1/2}\rceil$ windows of length
$h_{T}=(\log T)^{-1/2}$. The average number of on-line ordinates per
window is $\ge c'(\log T)^{1/2}$, so some window contains
$m(T)\ge c'(\log T)^{1/2}$ distinct ordinates
$\gamma$; distinct on-line zeros have distinct ordinates, and each
yields the zero $a=\tfrac14+i\gamma/2$ of $\Theta$ by
Lemma~\ref{lem:inner} (halving ordinates halves window lengths, which
we absorb into $h_{T}$). Let $y_{T}$ be the window center (after
halving) and $\varepsilon$ the offsets, so
$|\varepsilon|\le h_{T}/4\to0$. Finally $y_{T}\to\infty$: any bounded
region contains finitely many zeros, while $m(T)\to\infty$.
\end{proof}

\begin{lemma}\label{lem:smallTheta}
Fix $\delta\in(0,\tfrac14)$ and let $\mu_{n}=\delta+iy_{n}$ with
$y_{n},m_{n},\varepsilon_{n,j}$ as in Lemma~\ref{lem:clusters}. Then
\[
  |\Theta(\mu_{n})|
  \;\le\;
  q_{n,\delta}^{\,m_{n}}
  \longrightarrow 0
  \qquad (n\to\infty),
  \qquad
  q_{n,\delta}^{2}
  =\frac{(\tfrac14-\delta)^{2}+\max_{j}\varepsilon_{n,j}^{2}}
        {(\tfrac14+\delta)^{2}}
  \;\xrightarrow[n\to\infty]{}\;
  \Bigl(\frac{\tfrac14-\delta}{\tfrac14+\delta}\Bigr)^{2}<1 .
\]
\end{lemma}

\begin{proof}
Let $B_{n}$ be the finite Blaschke product over the cluster zeros
$a_{n,j}$ (with $b_{a}(s)=\frac{s-a}{s+\bar a}$). Since these are zeros
of the inner function $\Theta$, we have $\Theta=B_{n}\Theta_{n}$ with
$\Theta_{n}$ inner, so $|\Theta(\mu_{n})|\le|B_{n}(\mu_{n})|$. For each
factor,
\[
  |b_{a_{n,j}}(\mu_{n})|^{2}
  =\frac{(\tfrac14-\delta)^{2}+\varepsilon_{n,j}^{2}}
        {(\tfrac14+\delta)^{2}+\varepsilon_{n,j}^{2}}
  \le q_{n,\delta}^{2},
\]
and there are $m_{n}$ factors. Since $q_{n,\delta}$ is eventually
bounded by some $q_{\delta}<1$ and $m_{n}\to\infty$, the product tends
to $0$.
\end{proof}

\section{Main theorem}

\begin{theorem}\label{thm:main}
Unconditionally, $\|Z(t)\|=1$ for every $t\ge0$. Equivalently, the
growth bound of the semigroup $Z$ is
$\omega_{0}=\lim_{t\to\infty}t^{-1}\log\|Z(t)\|=0$, and
\[
  \dist_{H^{\infty}}\bigl(e^{-ts},\,\Theta\Hinf\bigr)=1
  \qquad\text{for every } t>0 .
\]
\end{theorem}

\begin{proof}
Fix $t>0$ and $\delta\in(0,\tfrac14)$, and let $\mu=\mu_{n}$ as in
Lemma~\ref{lem:smallTheta}, writing $\theta=\Theta(\mu)$. Using
$k^{\Theta}_{\mu}=k_{\mu}-\bar\theta\,\Theta k_{\mu}$ and the facts
(i) of Section~2,
\[
  Z(t)^{*}k^{\Theta}_{\mu}
  = T_{\bar\varphi_{t}}k_{\mu}
    -\bar\theta\,T_{\bar\varphi_{t}}(\Theta k_{\mu})
  = e^{-t\bar\mu}\,k_{\mu}
    -\bar\theta\,T_{\bar\varphi_{t}}(\Theta k_{\mu}),
\]
where $\|T_{\bar\varphi_{t}}(\Theta k_{\mu})\|\le
\|\varphi_{t}\|_{\infty}\|\Theta k_{\mu}\|=\|k_{\mu}\|$ because
$\|\varphi_{t}\|_{\infty}=1$ and multiplication by the inner function
$\Theta$ is isometric on $\HH$. Hence
\[
  \frac{\|Z(t)^{*}k^{\Theta}_{\mu}\|}{\|k^{\Theta}_{\mu}\|}
  \;\ge\;
  \frac{e^{-t\delta}-|\theta|}{\sqrt{1-|\theta|^{2}}}\,
  \qquad\text{(using } \|k^{\Theta}_{\mu}\|
    =\sqrt{1-|\theta|^{2}}\,\|k_{\mu}\| \text{)} .
\]
By Lemma~\ref{lem:smallTheta}, $|\theta|=|\Theta(\mu_{n})|\to0$ as
$n\to\infty$, so $\|Z(t)\|=\|Z(t)^{*}\|\ge e^{-t\delta}$. Letting
$\delta\downarrow0$ gives $\|Z(t)\|\ge1$; the reverse inequality holds
since $Z(t)$ is a contraction. For a contraction semigroup,
$\|Z(t_{0})\|<1$ for a single $t_{0}$ would force
$\omega_{0}\le t_{0}^{-1}\log\|Z(t_{0})\|<0$ by submultiplicativity, so
$\|Z(t)\|\equiv1$ is equivalent to $\omega_{0}=0$. The distance
statement follows from Sarason's formula (Section~2, (ii)).
\end{proof}

\begin{corollary}\label{cor:gap}
Assume RH. Then every eigenvalue of the generator $A$ has real part
$-\tfrac14$, while $\omega_{0}(Z)=0$: the spectral bound--growth bound
gap equals $\tfrac14$, i.e.\ RH does \emph{not} imply uniform
exponential stability of the scalar model semigroup.
\end{corollary}

\begin{corollary}\label{cor:LP}
Let $\widetilde Z(t)$ be any $C_{0}$-semigroup similar to $Z(t)$ with
similarity constants $c\|f\|\le\|Vf\|\le C\|f\|$ (in particular, any
unitarily equivalent realization, such as the identification of the
Lax--Phillips energy semigroup with the scalar compression proposed in
\cite{Uetake}). Then $\omega_{0}(\widetilde Z)=0$; in particular the
operator-norm bound \eqref{eq:naive} fails for $\widetilde Z$,
unconditionally. This does not contradict the Lax--Phillips criterion
in its per-vector form: for each zero $a$ of $\Theta$ the eigenvector
$k_{a}$ satisfies $\|Z(t)^{*}k_{a}\|=e^{-t\re a}\|k_{a}\|$, with
$\re a=\tfrac14$ under RH, and $Z(t)f\to0$ strongly for every
$f\in\KT$ since $\Theta$ is inner (the model semigroup is of class
$C_{\cdot0}$).
\end{corollary}

\begin{remark}\label{rem:extra-factor}
The full scattering matrix $S_{c}(s)=-\frac{s-1/2}{s+1/2}\Theta(s)$ is
inner with the same zeros as $\Theta$ plus a simple zero at
$s=\tfrac12$. Lemmas~\ref{lem:clusters}--\ref{lem:smallTheta} and the
proof of Theorem~\ref{thm:main} apply verbatim, so the conclusion holds
for $K_{S_{c}}$ as well.
\end{remark}

\begin{remark}
The proof exhibits the mechanism concretely: normalized reproducing
kernels $k^{\Theta}_{\mu_{n}}/\|k^{\Theta}_{\mu_{n}}\|$ at height
$y_{n}$ and abscissa $\delta$ are asymptotic eigenvectors of
$Z(t)^{*}$ with asymptotic eigenvalue modulus $e^{-t\delta}$. Clusters
of $m_{n}$ zeros violate the Carleson separation condition with
unbounded multiplicity, so the eigenvectors of $Z(t)^{*}$ do not form a
Riesz basis and no uniform decay rate can be assembled from the
eigenvalues. Equivalently, $|\Theta|$ fails to be bounded below on
every strip $0<\re s<\delta$, so the resolvent of $A$ is unbounded on
every vertical line $\re\lambda=-\delta$, and $\omega_{0}=0$ also
follows from the Gearhart--Pr\"uss theorem \cite{GearhartPruess}. In
the framework of \cite{NabokoRomanov}, the semigroup has a spectral
singularity at infinity.
\end{remark}

\begin{remark}
Only the existence of $\gg T\log T$ zeros on the critical line is used,
so Selberg's theorem suffices; neither simplicity (Levinson, Conrey)
nor any information about off-line zeros enters. In particular the
argument is identical whether RH holds or fails, which is why it can
refute only the operator-norm \emph{formulation} \eqref{eq:naive} and
carries no information about RH itself.
\end{remark}

\begin{thebibliography}{99}

\bibitem{Selberg} A. Selberg,
\emph{On the zeros of Riemann's zeta-function},
Skr. Norske Vid. Akad. Oslo I (1942), no. 10, 1--59.

\bibitem{Sarason} D. Sarason,
\emph{Generalized interpolation in $H^{\infty}$},
Trans. Amer. Math. Soc. \textbf{127} (1967), 179--203.

\bibitem{LaxPhillips} P. D. Lax and R. S. Phillips,
\emph{Scattering Theory for Automorphic Functions},
Annals of Mathematics Studies 87, Princeton Univ. Press, 1976.

\bibitem{FaddeevPavlov} L. D. Faddeev and B. S. Pavlov,
\emph{Scattering theory and automorphic functions},
J. Soviet Math. \textbf{3} (1975), 522--548.

\bibitem{LaxTodo} P. D. Lax,
\emph{[TODO: exact source of equation (54) --- verify citation and
whether the equation is stated in operator norm or per vector]}.

\bibitem{Uetake} Y. Uetake,
\emph{[TODO: exact title, journal, year 2007 --- the unitary
identification of the energy semigroup with the scalar compression]}.

\bibitem{GMR} S. R. Garcia, J. Mashreghi, and W. T. Ross,
\emph{Introduction to Model Spaces and their Operators},
Cambridge Univ. Press, 2016.

\bibitem{GearhartPruess} J. Pr\"uss,
\emph{On the spectrum of $C_{0}$-semigroups},
Trans. Amer. Math. Soc. \textbf{284} (1984), 847--857.

\bibitem{EngelNagel} K.-J. Engel and R. Nagel,
\emph{One-Parameter Semigroups for Linear Evolution Equations},
Springer GTM 194, 2000.

\bibitem{NabokoRomanov} S. Naboko and R. Romanov,
\emph{Spectral singularities and asymptotics of contractive
semigroups. I},
Acta Sci. Math. (Szeged) \textbf{70} (2004), 379--403.
[TODO: verify volume/pages]

\bibitem{Titchmarsh} E. C. Titchmarsh,
\emph{The Theory of the Riemann Zeta-Function}, 2nd ed.,
Oxford Univ. Press, 1986.

\end{thebibliography}

\end{document}
